% =============================================================================
%  proof_complete_standalone.tex
%
%  Self-contained version of proof_complete.tex.  It reuses the same preamble
%  (includes.tex / MacrosAndFigures.tex / Bib.bib) as main.tex.
%
%  Layout:
%    1. Background        -- general facts on symmetric powers of torsors
%                            (over an abstract base) that the proof quotes;
%    2. Geometric setting -- specialise to a curve: bounded ramification,
%                            the symmetric tower over U^(d), compactification,
%                            blowup;
%    3. The master diagram;
%    4. The unramifiedness theorem and its proof.
%
%  Compile with:  latexmk -pdf proof_complete_standalone.tex   (uses biber)
% =============================================================================
\documentclass[11pt]{article}
\usepackage{geometry}\geometry{margin=1in}
\usepackage[T1]{fontenc}
\usepackage{ifthen}

\setlength{\parindent}{1em}
\setlength{\parskip}{0.5em}

\input{includes}
\input{MacrosAndFigures}

\usepackage{csquotes}
\usepackage[
backend=biber,
style=alphabetic,
sorting=ynt
]{biblatex}
\addbibresource{Bib.bib}

\newboolean{ThesisIntro}
\setboolean{ThesisIntro}{false}

\title{Unramifiedness of $\cP^{[d]_G}$ at the exceptional divisor\\[2pt]
       {\large (standalone excerpt)}}
\author{Assaf Marzan}
\date{\today}

\begin{document}
\maketitle

\noindent
\textbf{Roadmap.} \Cref{sec:background} collects, without re-deriving them, the
facts about symmetric powers of torsors that the proof uses, stated over an
abstract base.  \Cref{sec:setting} specialises this to a smooth projective
curve: it recalls the notion of bounded ramification and fixes the geometric
data (the symmetric tower over $U^{(d)}$, the compactification $C' \to C$, the
blowups at the points at infinity).  \Cref{sec:master} assembles everything into
a single master diagram, and \Cref{sec:main} states and proves the
unramifiedness theorem.

% #############################################################################
\section{Background}\label{sec:background}
% #############################################################################

Throughout this section $S$ is a base scheme, $G$ a finite abelian group,
$H \subseteq G$ a subgroup, and $d \geq 0$ an integer.  All bases are assumed
Zariski-locally quasi-projective over $S$, so that the symmetric-power quotients
below exist as schemes.

%------------------------------------------------------------------------------
\subsection{Symmetric powers of torsors}\label{subsec:sym-powers}
%------------------------------------------------------------------------------

Let $\cP$ be a $G$-torsor over an $S$-scheme $X$.  For each
$i \in \llbracket 1,d\rrbracket$ write $p_i : X^{\times_S d} \to X$ for the
$i$-th projection and consider the $G$-torsor
\[
p_1^{-1}\cP \otimes^G \cdots \otimes^G p_d^{-1}\cP \;=\; \KG \backslash \cP^{\times_S d}
\]
over $X^{\times_S d}$, where $\KG \subseteq G^d$ is the kernel of the
multiplication morphism $G^d \to G$ (the dependence on $d$ is suppressed in the
notation).  The symmetric group $S_d$ acts on the right on
$\KG \backslash \cP^{\times_S d}$ by $(p_i)_i \cdot \sigma = (p_{\sigma(i)})_i$,
and this action commutes with the left $G$-action.

\begin{prop}[\cite{Guignard2018}, Proposition 2.32]\label{prop:symmetric_power_torsor}
The right action of $S_d$ on $\KG \backslash \cP^{\times_S d}$ is admissible, so
that the quotient $\cP^{[d]}$ of $\KG \backslash \cP^{\times_S d}$ by $S_d$
exists as a scheme over $S$.  Moreover, the canonical morphism
$\cP^{[d]} \to \mathrm{Sym}_S^d(X)$ is a $G$-torsor, and the morphism
\[
p_1^{-1}\cP \otimes^G \cdots \otimes^G p_d^{-1}\cP \to r^{-1}\cP^{[d]},
\]
where $r : X^{\times_S d} \to \mathrm{Sym}_S^d(X)$ is the canonical projection,
is an isomorphism of $G$-torsors over $X^{\times_S d}$.
\end{prop}

When $G$ is abelian the condition $\prod_{i=1}^d g_i = 1$ defining $\KG$ is
invariant under permuting the coordinates, so the (left) $S_d$-action normalises
the $\KG$-action.  Writing $\Xi = \KG \rtimes S_d$ for the resulting finite
group, acting on $\cP^{\times_S d}$ by
\[
\bigl((g_1,\dots,g_d),\sigma\bigr) \star (p_1,\dots,p_d)
   = \bigl(g_1 p_{\sigma^{-1}(1)},\,\dots,\,g_d p_{\sigma^{-1}(d)}\bigr),
\]
one has $(\KG \backslash \cP^{\times_S d})/S_d \cong \Xi \backslash \cP^{\times_S d}$.

\begin{cor}\label{cor:sym-power-quotient}
There is a natural canonical morphism
\[
q : S_d \backslash \cP^{\times_S d} \longrightarrow \Xi \backslash \cP^{\times_S d}
\qquad\text{equivalently}\qquad
q : S_d \backslash \cP^{\times_S d} \longrightarrow (\KG \backslash \cP^{\times_S d})/S_d ,
\]
the comparison between the symmetric power of $\cP$ \emph{as a scheme} and the
symmetric power of $\cP$ \emph{as a $G$-torsor}.  We write $\cP^{(d)}$ for the
former and $\cP^{[d]}$ for the latter, so that $q : \cP^{(d)} \to \cP^{[d]}$.
\end{cor}
\begin{proof}
$S_d$ is a subgroup of $\Xi$, hence any $\Xi$-invariant morphism is
automatically $S_d$-invariant.
\end{proof}

%------------------------------------------------------------------------------
\subsection{The canonical decomposition of a torsor}\label{subsec:decomposition}
%------------------------------------------------------------------------------

\begin{prop}\label{prop:quotient_torsors_finite}\label{prop:quotient_torsors}
Let $\cP \to X$ be a $G$-torsor in $X_{\text{\'et}}$ and let $H \subseteq G$ be
a subgroup, with quotient map $\pi : G \to G/H$.  Then $\cP \to X$ factors
\[
\cP \xrightarrow{\;\phi\;} \cP_{G/H} \xrightarrow{\;\psi\;} X,
\]
uniquely up to a unique isomorphism commuting with the projections, where
\begin{enumerate}
    \item $\cP_{G/H} := \cP \times^G (G/H)$ is a $G/H$-torsor over $X$, and
    \item $\cP$ is an $H$-torsor over the scheme $\cP_{G/H}$ via $\phi$.
\end{enumerate}
\end{prop}

%------------------------------------------------------------------------------
\subsection{The tower of symmetric quotients}\label{subsection:quotient_torsors_towers}
%------------------------------------------------------------------------------

Fix a $G$-torsor $\cP \to X$ and a subgroup $H \subseteq G$, with canonical
factorisation $\cP \xrightarrow{\phi} \cP_{G/H} \xrightarrow{\psi} X$ as in
\Cref{prop:quotient_torsors_finite} ($\phi$ an $H$-torsor, $\psi$ a
$G/H$-torsor).  We record how the symmetric-power construction of
\Cref{subsec:sym-powers} interacts with this factorisation.

\paragraph{Summation kernels.}
Besides $\KG = \ker(G^d \to G)$ we need
\[
K_H := \ker(H^d \to H), \qquad K_{G/H} := \ker\bigl((G/H)^d \to G/H\bigr).
\]
The snake lemma applied to
\[
\begin{tikzcd}[row sep=large, column sep=large]
0 \arrow[r] & H^d \arrow[r] \arrow[d, twoheadrightarrow] & G^d \arrow[r] \arrow[d, twoheadrightarrow] & (G/H)^d \arrow[r] \arrow[d, twoheadrightarrow] & 0 \\
0 \arrow[r] & H \arrow[r] & G \arrow[r] & G/H \arrow[r] & 0,
\end{tikzcd}
\]
(the vertical maps being the surjective summation morphisms) gives a short exact
sequence
\begin{equation}\label{eq:summation_kernels_ses}
0 \to K_H \to \KG \to K_{G/H} \to 0,
\end{equation}
hence $\KG/K_H \cong K_{G/H}$.  Inside $G^d$ we have the chains of subgroups
\[
K_H \;\subseteq\; \KG \;\subseteq\; H^d + \KG, \qquad H^d \;\subseteq\; H^d + \KG .
\]

\paragraph{Four quotient schemes.}
Let $S_d$ act on $\cP^{\times_S d}$ as in \Cref{cor:sym-power-quotient}, and let
each of the subgroups above act diagonally via the $G$-action on $\cP$.  Set
\[
\begin{aligned}
    \cP^{[d]_G}        &:= (\KG \rtimes S_d)\backslash \cP^{\times_S d},
        &&\text{the symmetric $G$-power of $\cP$;}\\
    \cP^{[d]_H}        &:= (K_H \rtimes S_d)\backslash \cP^{\times_S d},
        &&\text{the symmetric $H$-power of $\cP$ relative to $\phi$;}\\
    (\cP_{G/H})^{(d)}  &:= (H^d \rtimes S_d)\backslash \cP^{\times_S d},
        &&\text{the symmetric power of the scheme $\cP_{G/H}$;}\\
    (\cP_{G/H})^{[d]}  &:= \bigl((H^d + \KG)\rtimes S_d\bigr)\backslash \cP^{\times_S d},
        &&\text{the symmetric $G/H$-power of $\cP_{G/H}$.}
\end{aligned}
\]

\begin{prop}\label{prop:torsor_tower_diagram}
The four quotient schemes above exist and fit into a canonical commutative
diagram
\[
\begin{tikzcd}[row sep=large, column sep=large]
    \cP^{[d]_H} \arrow[r, "H"] \arrow[d] & (\cP_{G/H})^{(d)} \arrow[d, "q"] \\
    \cP^{[d]_G} \arrow[r, "H"'] \arrow[dr, "G"', bend right=20] & (\cP_{G/H})^{[d]} \arrow[d, "G/H"] \\
    & X^{(d)},
\end{tikzcd}
\]
in which each labelled arrow is a torsor under the indicated constant group, the
right-hand vertical $q$ is the canonical projection of
\Cref{cor:sym-power-quotient} applied to $\cP_{G/H} \to X$, and the upper square
is Cartesian:
\[
\cP^{[d]_H} \cong \cP^{[d]_G} \times_{(\cP_{G/H})^{[d]}} (\cP_{G/H})^{(d)}.
\]
The unlabelled left-hand vertical $\cP^{[d]_H} \to \cP^{[d]_G}$ is then the base
change of $q$ along $\cP^{[d]_G} \to (\cP_{G/H})^{[d]}$.  Neither $q$ nor this
left-hand vertical is, in general, a torsor under a constant group scheme.
\end{prop}
\begin{proof}
The four labelled arrows are torsors under the indicated constant groups by
applying \Cref{prop:symmetric_power_torsor} and
\Cref{prop:quotient_torsors_finite}:
\begin{itemize}
    \item to the $G$-torsor $\cP \to X$, giving the $G$-torsor
          $\cP^{[d]_G} \to X^{(d)}$ (the diagonal arrow);
    \item to the $H$-torsor $\phi : \cP \to \cP_{G/H}$ (with $\cP_{G/H}$ as
          base), giving the $H$-torsor $\cP^{[d]_H} \to (\cP_{G/H})^{(d)}$
          (the top arrow);
    \item to the $G/H$-torsor $\psi : \cP_{G/H} \to X$, giving the $G/H$-torsor
          $(\cP_{G/H})^{[d]} \to X^{(d)}$ (the bottom-right vertical) and, via
          \Cref{cor:sym-power-quotient}, the comparison map
          $q : (\cP_{G/H})^{(d)} \to (\cP_{G/H})^{[d]}$.
\end{itemize}
The factorisation $\cP^{[d]_G} \xrightarrow{H} (\cP_{G/H})^{[d]} \xrightarrow{G/H} X^{(d)}$
of the diagonal arrow is the canonical decomposition of
\Cref{prop:quotient_torsors_finite} applied to $\cP^{[d]_G} \to X^{(d)}$ and
$H \subseteq G$: both $\cP^{[d]_G} \times^G (G/H)$ and $(\cP_{G/H})^{[d]}$ are
realised as the admissible quotient
$\bigl((H^d + \KG)\rtimes S_d\bigr)\backslash \cP^{\times_S d}$.

For the Cartesian property, set
$F := \cP^{[d]_G} \times_{(\cP_{G/H})^{[d]}} (\cP_{G/H})^{(d)}$.  Base-changing
the $H$-torsor $\cP^{[d]_G} \to (\cP_{G/H})^{[d]}$ along $q$ exhibits $F$ as an
$H$-torsor over $(\cP_{G/H})^{(d)}$.  The universal property of the fibre
product gives a canonical $H$-equivariant morphism $\cP^{[d]_H} \to F$ over
$(\cP_{G/H})^{(d)}$; any morphism of $H$-torsors is an isomorphism, so
$\cP^{[d]_H} \cong F$ and the upper square is Cartesian.  Via this square the
left vertical $\cP^{[d]_H} \to \cP^{[d]_G}$ is identified with the base change
of $q$ along $\cP^{[d]_G} \to (\cP_{G/H})^{[d]}$.
\end{proof}

% #############################################################################
\section{The geometric setting}\label{sec:setting}
% #############################################################################

%------------------------------------------------------------------------------
\subsection{Bounded ramification along a curve}\label{subsec:bounded-ram}
%------------------------------------------------------------------------------

Let $C$ be a smooth projective geometrically connected curve over a field $k$,
let $\mdls$ be a fixed modulus, and put $U = C \setminus \mdls$.  Let
$\cP \to U$ be a $G$-torsor with $G$ finite abelian.  For a closed point $P$ of
$C$, passing to the completion $\widehat{\Oo}_{C,P}$ gives the complete local
field $\widehat{K}_P$, and restricting $\cP$ to $\Spec(\widehat{K}_P)$ yields a
$G$-torsor over a field, hence a disjoint union of spectra of finite separable
field extensions
\[
\cP|_{\Spec(\widehat{K}_P)} \cong \bigsqcup \Spec(F),
\]
in which the extensions $F/\widehat{K}_P$ are all isomorphic and Galois,
$H = \Gal(F/\widehat{K}_P)$ is the stabiliser of a connected component (a
subgroup of $G$), and there are $[G:H]$ components.  Write $H^r$ for the $r$-th
higher ramification group of $F/\widehat{K}_P$.

\begin{definition}\label{definition:local_ramification_boundness}
The extension $F/\widehat{K}_P$ \textit{has ramification bounded by $r$} if the
ramification group $H^r$ is trivial.  We say \textit{$\cP$ has ramification
bounded by $r$ at $P$} if the corresponding local extensions $F/\widehat{K}_P$
satisfy this condition.
\end{definition}

%------------------------------------------------------------------------------
\subsection{The symmetric tower over \texorpdfstring{$U^{(d)}$}{the symmetric power}}
\label{subsec:tower-over-U}
%------------------------------------------------------------------------------

For the rest of the note we fix, in addition to the data above:
\begin{itemize}
  \item a normal subgroup $H \trianglelefteq G$;
  \item a point $p_\infty \in \mdls$ and an integer $d \geq 1$;
\end{itemize}
and we assume that $\cP \to U$ has ramification bounded by $d$ at $p_\infty$ in
the sense of \Cref{definition:local_ramification_boundness} (i.e.\ the $d$-th
higher ramification group of the extension of DVRs at $p_\infty$ is trivial).

Apply the construction of \Cref{subsection:quotient_torsors_towers} with
$S = \Spec k$ and base $X = U$.  This produces the four quotient schemes
\[
\cP^{[d]_H},\qquad \cP^{[d]_G},\qquad (\cP_{G/H})^{(d)},\qquad (\cP_{G/H})^{[d]},
\]
and, by \Cref{prop:torsor_tower_diagram}, the canonical commutative diagram with
Cartesian upper square
\[
\begin{tikzcd}[row sep=large, column sep=large]
\cP^{[d]_H} \ar[r,"H"] \ar[d]
  & (\cP_{G/H})^{(d)} \ar[d,"q"] \\
\cP^{[d]_G} \ar[r,"H"'] \ar[dr,"G"',bend right=15]
  & (\cP_{G/H})^{[d]} \ar[d,"G/H"] \\
  & U^{(d)},
\end{tikzcd}
\]
in which the arrows labelled $H$, $G/H$ and $G$ are torsors under the indicated
finite constant groups (in particular, \'etale).

%------------------------------------------------------------------------------
\subsection{Compactification and blowup}\label{subsec:compactify-blowup}
%------------------------------------------------------------------------------

The $G/H$-torsor $\cP_{G/H} \to U$ compactifies to a finite morphism of smooth
projective curves $f \colon C' \to C$; fix $p'_\infty \in f^{-1}(p_\infty)$.
Taking $d$-fold symmetric powers gives the finite morphism
\[
f^{(d)} \colon C'^{(d)} \to C^{(d)}, \qquad d \cdot p'_\infty \longmapsto d \cdot p_\infty .
\]
Let
\[
\pi_C \colon X_C \to C^{(d)},\qquad \pi_{C'} \colon X_{C'} \to C'^{(d)}
\]
be the blowups at $d \cdot p_\infty$ and $d \cdot p'_\infty$ respectively, with
exceptional divisors $E_C, E_{C'}$ and generic points $\eta_C, \eta_{C'}$.  The
local rings at these generic points are the discrete valuation rings
\[
V_C := \mathcal{O}_{X_C,\eta_C} \subset K(C^{(d)}),\qquad
V_{C'} := \mathcal{O}_{X_{C'},\eta_{C'}} \subset K(C'^{(d)}).
\]

% #############################################################################
\section{The master diagram}\label{sec:master}
% #############################################################################

The torsor tower over $U^{(d)}$ embeds into the projective geometry of $C^{(d)}$
and $C'^{(d)}$, which in turn carry the canonical blowups at the points at
infinity.  Assembling all of this:
\[
\begin{tikzcd}[row sep=2.4em, column sep=2.2em]
\cP^{[d]_H} \ar[r,"H"] \ar[d]
  & (\cP_{G/H})^{(d)} \ar[d,"q"] \ar[r,hookrightarrow,"\iota'"]
  & C'^{(d)} \ar[dd,"f^{(d)}"]
  & X_{C'} \ar[l,"\pi_{C'}"'] \ar[dd,dashed,"f_X"] \\
\cP^{[d]_G} \ar[r,"H"'] \ar[dr,"G"',bend right=10]
  & (\cP_{G/H})^{[d]} \ar[d,"G/H"]
  & & \\
  & U^{(d)} \ar[r,hookrightarrow,"\iota"']
  & C^{(d)}
  & X_C \ar[l,"\pi_C"]
\end{tikzcd}
\]
Reading the diagram from left to right:
\begin{description}
  \item[Left block (columns 1--2):] the torsor tower of
        \Cref{prop:torsor_tower_diagram} over $U^{(d)}$ --- $H$-torsors
        horizontally, the $G/H$-torsor on the lower right, and the composite
        $G$-torsor $\cP^{[d]_G} \to U^{(d)}$.
  \item[Middle block (column 3):] the compactifications.  Here $\iota$ is the
        canonical open immersion $U^{(d)} \hookrightarrow C^{(d)}$, and $\iota'$
        is the open immersion $(\cP_{G/H})^{(d)} \hookrightarrow C'^{(d)}$
        coming from $\cP_{G/H} \subset C'$.  Restricting $f^{(d)}$ along
        $\iota'$ recovers the canonical map $(\cP_{G/H})^{(d)} \to U^{(d)}$,
        i.e.\ the composite of $q$ with the $G/H$-torsor.
  \item[Right block (column 4):] the blowups at the diagonal points at infinity.
        The dashed arrow $f_X \colon X_{C'} \dashrightarrow X_C$ is the rational
        lift of $f^{(d)}$, defined on the strict transform; on generic points of
        the exceptional divisors it sends $\eta_{C'} \mapsto \eta_C$, inducing
        an extension of DVRs $V_C \hookrightarrow V_{C'}$.
\end{description}

% #############################################################################
\section{Unramifiedness of \texorpdfstring{$\cP^{[d]_G}$}{P\^{}[d]\_G} at \texorpdfstring{$\eta_C$}{eta\_C}}\label{sec:main}
% #############################################################################

\begin{theorem}\label{thm:main}
Suppose
\begin{enumerate}
  \item[\textnormal{(H1)}] $(\cP_{G/H})^{[d]} \to U^{(d)}$ is unramified at
        $\eta_C$; and
  \item[\textnormal{(H2)}] $\cP^{[d]_H} \to (\cP_{G/H})^{(d)}$ is unramified at
        $\eta_{C'}$.
  \item $G=Z/p^rZ$ is cyclic of prime power order.
\end{enumerate}
Then $\cP^{[d]_G} \to U^{(d)}$ is unramified at $\eta_C$.
\end{theorem}

\begin{proof}
We may assume that $\cP \to U$ is connected and that the monodromy
homomorphism $\pi_1(U^{(d)}) \to G$ associated to $\cP^{[d]_G}$ is surjective:
otherwise replace $G$ by its image and $H$ by the intersection, both of which
inherit the hypotheses. We may also assume $H \neq 0$, since otherwise
$(\cP_{G/H})^{[d]} = \cP^{[d]_G}$ and (H1) coincides with the conclusion.

Under these reductions, both function fields
\[
E := K(\cP^{[d]_G}), \qquad F := K\bigl((\cP_{G/H})^{[d]}\bigr)
\]
are single fields, Galois over $K_C := \Frac V_C$ with groups $G$ and $G/H$
respectively, and the canonical factorisation of
\Cref{prop:torsor_tower_diagram} furnishes a tower
$K_C \subseteq F \subseteq E$ with $\Gal(E/F) = H$.

\paragraph{Step 1: A compatible place of $F$ above $V_C$.}
Since $F \subseteq K_{C'} := \Frac V_{C'}$, the DVR $V_{C'}$ restricts to a DVR
$V_F$ on $F$, and $V_F$ lies above $V_C$. By (H1), $V_F/V_C$ is an unramified
extension of DVRs. Fix a place $v_E$ of $E$ above $V_F$ and a place $v_\star$
of the compositum $E \cdot K_{C'}$ (formed in a fixed algebraic closure of
$K_C$) above both $v_E$ and $V_{C'}$. Towers of unramified DVR extensions
being unramified, it suffices to prove that the inertia
\[
I \;:=\; I_{v_E}(E/F) \;\leq\; \Gal(E/F) \;=\; H
\]
is trivial.

\paragraph{Step 2: $E$ does not embed into $K_{C'}$.}
Let $L := K(\cP^{\times d})$. Since $\cP \to U$ is a connected $G$-torsor,
$\cP^{\times d} \to U^d$ is a connected $G^d$-torsor, and the natural
$S_d$-action on $\cP^{\times d}$ covering the $S_d$-action on $U^d$ exhibits
$L/K_C$ as a Galois extension with group
\[
\widetilde G \;:=\; G^d \rtimes S_d .
\]
By the description of $\cP^{[d]_G}$ as $(\KG \rtimes S_d)\backslash \cP^{\times d}$
and the fact that $\Xi := \KG \rtimes S_d$ is normal in $\widetilde G$ with
quotient $\widetilde G/\Xi = G^d/\KG = G$, Galois correspondence gives
\[
E \;=\; L^{\Xi} .
\]
On the other hand, $\cP \to \cP_{G/H}$ is an $H$-torsor, so
$\cP^{\times d}/H^d = (\cP_{G/H})^{\times d}$, and consequently
\[
K_{C'} \;=\; K\!\bigl((\cP_{G/H})^{(d)}\bigr)
       \;=\; \bigl(L^{H^d}\bigr)^{S_d}
       \;=\; L^{H^d \rtimes S_d} .
\]
Galois correspondence in $L/K_C$ now reads
\[
E \subseteq K_{C'}
   \;\iff\; \Xi \supseteq H^d \rtimes S_d
   \;\iff\; \KG \supseteq H^d .
\]
For any $h \in H$, the tuple $(h,0,\dots,0) \in H^d$ has summation $h$ in $G$;
hence $\KG \supseteq H^d$ forces $h = 0$ for every $h \in H$, contradicting
$H \neq 0$. Therefore
\[
E \;\not\subseteq\; K_{C'} .
\]

\paragraph{Step 3: Inertia comparison via (H2).}
Set $M := E \cap K_{C'}$, so that $F \subseteq M \subsetneq E$ by Step 2. The
compositum $E \cdot K_{C'}$ is Galois over $K_{C'}$ with group
\[
N \;:=\; \Gal\bigl(E \cdot K_{C'}/K_{C'}\bigr) \;=\; \Gal(E/M)
     \;\leq\; \Gal(E/F) \;=\; H,
\]
where the inclusion $N \hookrightarrow H$ is restriction to $E$, and $N \neq 1$.

The Cartesian property of \Cref{prop:torsor_tower_diagram} gives
\[
\cP^{[d]_H} \;=\; \cP^{[d]_G} \times_{(\cP_{G/H})^{[d]}} (\cP_{G/H})^{(d)},
\]
so the generic fibre of $\cP^{[d]_H}$ over $\Spec K_{C'}$ is
$E \otimes_F K_{C'}$, whose connected components are the conjugates of
$E \cdot K_{C'}$ under $\Gal(F/K_C) = G/H$. By (H2), each component is
unramified at $V_{C'}$; in particular,
\[
I' \;:=\; I_{v_\star}\!\bigl(E\cdot K_{C'}/K_{C'}\bigr) \;=\; 1 .
\]
The standard formula for inertia in a compositum yields, as subgroups of $H$
via $N \hookrightarrow H$,
\[
I' \;=\; I \cap N .
\]
(One direction is immediate: $\sigma \in I'$ acts trivially on $\kappa(v_\star)$,
hence on $\kappa(v_E) \subseteq \kappa(v_\star)$, so $\sigma|_E \in I$. The
reverse uses $\kappa(v_\star) = \kappa(v_E)\cdot\kappa(V_{C'})$, valid because
the residue extension $\kappa(v_E)/\kappa(V_F)$ is separable: it is acted on
faithfully by the cyclic $p$-group $\Gal(E/F)/I$, which has no purely
inseparable target.)

Hence $I \cap N = 1$ inside $H$.

\paragraph{Step 4: Cyclic prime-power structure forces $I = 1$.}
Since $G = \Z/p^r\Z$ is cyclic of prime-power order, so is $H$. The subgroups
of $H$ form a chain under inclusion; consequently, two subgroups $A, B \leq H$
satisfy $A \cap B = 1$ if and only if $A = 1$ or $B = 1$. Applied to
$I \cap N = 1$ with $N \neq 1$, we obtain $I = 1$.

Thus $E/F$ is unramified at $V_F$. Combined with the unramified extension
$V_F/V_C$ of Step~1, this proves that $E/K_C$ is unramified at $V_C$, i.e.\
that $\cP^{[d]_G} \to U^{(d)}$ is unramified at $\eta_C$.
\end{proof}

\begin{remark}
The cyclic prime-power hypothesis on $G$ enters \emph{only} in Step~4, where
it is genuinely needed: for a non-cyclic group such as $(\Z/2)^2$, two
non-trivial subgroups can intersect trivially, and $I \cap N = 1$ no longer
forces $I = 1$ given $N \neq 1$.

The standing assumption that $\cP$ has ramification bounded by $d$ at
$p_\infty$ is not used in the argument above; its role is upstream, ensuring
that the rational map $f_X \colon X_{C'} \dashrightarrow X_C$ sends
$\eta_{C'}$ to $\eta_C$ and so produces the finite extension of DVRs
$V_C \hookrightarrow V_{C'}$ on which Steps~1 and~3 rely.
\end{remark}


% #############################################################################
\section{A revised proof of \texorpdfstring{\Cref{thm:main}}{the unramifiedness theorem}}\label{sec:main-revised}
% #############################################################################

The proof of \Cref{thm:main} above suffers from (i) a labelling slip on
hypothesis~(H3), (ii) unjustified preliminary reductions, (iii) an incorrect
description of the connected components of $E \otimes_F K_{C'}$ in Step~3,
and (iv) an incomplete justification of the residue-field identity
$\kappa(v_\star) = \kappa(v_E) \cdot \kappa(V_{C'})$ underlying the inertia
comparison.  We restate the theorem with an explicit separability hypothesis
and give a tightened proof; the four-step skeleton is unchanged.

\begin{theorem}\label{thm:main-revised}
Assume
\begin{enumerate}
  \item[\textnormal{(H1)}] $(\cP_{G/H})^{[d]} \to U^{(d)}$ is unramified at
        $\eta_C$;
  \item[\textnormal{(H2)}] $\cP^{[d]_H} \to (\cP_{G/H})^{(d)}$ is unramified at
        $\eta_{C'}$;
  \item[\textnormal{(H3)}] $G = \Z/p^r\Z$ is cyclic of prime-power order;
  \item[\textnormal{(H4)}] the residue extension
        $\kappa(\eta_{C'})/\kappa(\eta_C)$ is separable.
\end{enumerate}
Then $\cP^{[d]_G} \to U^{(d)}$ is unramified at $\eta_C$.
\end{theorem}

Condition (H4) is automatic when $k$ has characteristic $0$, and more
generally when $\kappa(\eta_C)$ is perfect or when $V_{C'}/V_C$ is tame; in
those cases it imposes nothing new.

\begin{proof}
\paragraph{Preliminary reductions.}
We perform two reductions.
\begin{enumerate}
\item[\textnormal{(R1)}] \textit{$\cP \to U$ is connected.}  Let
      $\rho \colon \pi_1(U) \to G$ be the monodromy of $\cP$ and put
      $G_0 := \mathrm{Im}(\rho)$, $H_0 := H \cap G_0$.  Pick a connected
      component $\cP_0 \subseteq \cP$; it is a connected $G_0$-torsor with
      subgroup $H_0$, and $\cP$ is the disjoint union of $[G : G_0]$ copies
      of $\cP_0$ permuted by $G$.  Each of the four schemes in
      \Cref{prop:torsor_tower_diagram} for $(\cP, G, H)$ decomposes as a
      disjoint union of copies of the corresponding scheme for
      $(\cP_0, G_0, H_0)$, since the symmetric-power quotients commute with
      such disjoint unions.  Unramifiedness at $\eta_C$ (resp.\ $\eta_{C'}$)
      is detected componentwise, so (H1) and (H2) are inherited by
      $(\cP_0, G_0, H_0)$; (H3) is inherited because subgroups of cyclic
      $p$-groups are cyclic $p$-groups; (H4) is unchanged.  Replacing
      $(\cP, G, H)$ by $(\cP_0, G_0, H_0)$ we may assume $\cP$ is connected
      and $\rho$ is surjective.
\item[\textnormal{(R2)}] \textit{$H \neq 0$.}  If $H = 0$ then $\cP_{G/H} = \cP$ and
      $(\cP_{G/H})^{[d]} = \cP^{[d]_G}$; (H1) is then literally the
      conclusion.
\end{enumerate}
Under (R1)--(R2) the function fields
\[
E := K(\cP^{[d]_G}), \qquad F := K\bigl((\cP_{G/H})^{[d]}\bigr)
\]
are single fields, Galois over $K_C := \Frac V_C$ with groups $G$ and $G/H$
respectively, and \Cref{prop:torsor_tower_diagram} yields a tower
$K_C \subseteq F \subseteq E$ with $\Gal(E/F) = H$.

\paragraph{Step 1: A place $V_F$ above $V_C$, unramified by (H1).}
The comparison map $q : (\cP_{G/H})^{(d)} \to (\cP_{G/H})^{[d]}$ is generically
finite and dominant; combined with the open immersions
$(\cP_{G/H})^{(d)} \hookrightarrow C'^{(d)}$ and
$(\cP_{G/H})^{[d]} \to C^{(d)}$ (via the $G/H$-torsor and the open immersion
$U^{(d)} \hookrightarrow C^{(d)}$), this gives $F \subseteq K_{C'}$.
Restriction of $V_{C'}$ to $F$ defines a DVR $V_F$ with $V_F \mid V_C$.
Hypothesis (H1) asserts that \emph{every} place of $F$ above $V_C$ is
unramified; in particular $V_F/V_C$ is unramified.

Fix a place $v_E$ of $E$ above $V_F$ and, inside a chosen algebraic closure
of $K_C$, a place $v_\star$ of $E \cdot K_{C'}$ above both $v_E$ and
$V_{C'}$.  Since $H = \Gal(E/F)$ acts transitively on places of $E$ above
$V_F$, triviality of $I_{v_E}(E/F)$ implies triviality of inertia at every
place of $E$ above $V_F$; combined with the unramified tower step
$V_F/V_C$, that yields the conclusion.  It therefore suffices to prove
\[
I := I_{v_E}(E/F) \;=\; 1 \qquad \text{inside } H .
\]

\paragraph{Step 2: $E \not\subseteq K_{C'}$.}
Let $L := K(\cP^{\times d})$.  Since $\cP \to U$ is a connected $G$-torsor by
(R1), $\cP^{\times d} \to U^d$ is a connected $G^d$-torsor; the $S_d$-action
on $\cP^{\times d}$ covering the $S_d$-action on $U^d$ exhibits $L$ as a
Galois extension of $K_C$ with group
\[
\widetilde G := G^d \rtimes S_d .
\]
(Generically $L/K(U^d)$ is Galois with group $G^d$, $K(U^d)/K_C$ is Galois
with group $S_d$, and the two pieces fit into a semidirect product because
the $S_d$-permutation of factors normalises the $G^d$-action.)
By the definition of $\cP^{[d]_G}$ as $\Xi \backslash \cP^{\times d}$ with
$\Xi := \KG \rtimes S_d \trianglelefteq \widetilde G$ and quotient
$\widetilde G/\Xi = G^d/\KG = G$, Galois correspondence gives
$E = L^\Xi$.  Likewise
$(\cP_{G/H})^{(d)} = (H^d \rtimes S_d)\backslash \cP^{\times d}$ yields
$K_{C'} = L^{H^d \rtimes S_d}$.  Hence
\[
E \subseteq K_{C'}
   \iff \Xi \supseteq H^d \rtimes S_d
   \iff \KG \supseteq H^d .
\]
For any $h \in H$ the element $(h, 0, \dots, 0) \in H^d$ has summation $h$
in $G$, so $\KG \supseteq H^d$ would force $h = 0$ for every $h \in H$,
contradicting (R2).  Therefore $E \not\subseteq K_{C'}$.

\paragraph{Step 3: Inertia comparison via (H2).}
Set $M := E \cap K_{C'}$, so $F \subseteq M \subsetneq E$ (the inclusion
$F \subseteq M$ because $F \subseteq E \cap K_{C'}$; the proper inclusion
by Step~2).  Then $E \cdot K_{C'}$ is Galois over $K_{C'}$, and restriction
to $E$ identifies
\[
N := \Gal(E \cdot K_{C'} / K_{C'}) \;\cong\; \Gal(E/M)
    \;\leq\; \Gal(E/F) \;=\; H ,
\qquad N \neq 1 .
\]
The Cartesian property of \Cref{prop:torsor_tower_diagram} gives
\[
\cP^{[d]_H} \;=\; \cP^{[d]_G} \times_{(\cP_{G/H})^{[d]}} (\cP_{G/H})^{(d)} ,
\]
so the generic fibre of $\cP^{[d]_H} \to (\cP_{G/H})^{(d)}$ over $\Spec K_{C'}$
is $E \otimes_F K_{C'}$.  Since $E/F$ is Galois with group $H$, this is a
finite product of fields, each $K_{C'}$-isomorphic to $E \cdot K_{C'}$ and
Galois over $K_{C'}$ with group $N$; the number of factors is $[H : N]$
(equivalently, $|H|/|N|$).  By (H2) each factor is unramified at $V_{C'}$;
in particular
\[
I' \;:=\; I_{v_\star}\!\bigl(E \cdot K_{C'} / K_{C'}\bigr) \;=\; 1 .
\]

We claim that, under $N \hookrightarrow H$,
\begin{equation}\label{eq:inertia-intersection}
I' \;=\; I \cap N \qquad \text{inside } H .
\end{equation}
This is the standard \emph{inertia-in-compositum} identity for the pair of
Galois extensions $E/F$ and $K_{C'}/F$ inside their compositum, and holds
once the residue extension $\kappa(v_E)/\kappa(V_F)$ is separable
(cf.\ Bourbaki, \emph{Algèbre Commutative}, Ch.~VI, \S8; Serre,
\emph{Local Fields}, Ch.~I, \S7).
We verify separability from (H4):
$\kappa(V_F) = \kappa(\eta_C)$ since $V_F/V_C$ is unramified; the residue
extension $\kappa(v_\star)/\kappa(V_{C'})$ is separable because
$E \cdot K_{C'}/K_{C'}$ is unramified at $v_\star$ by (H2); and (H4) makes
$\kappa(V_{C'}) = \kappa(\eta_{C'})$ separable over $\kappa(\eta_C)$.
The subextension $\kappa(v_E)/\kappa(V_F) \hookrightarrow
\kappa(v_\star)/\kappa(\eta_C)$ is therefore contained in a separable
extension and is itself separable.

Both inclusions of \eqref{inertia-intersection} are now standard:
\begin{itemize}
\item If $\sigma \in I'$ then $\sigma$ acts trivially on $\kappa(v_\star)$,
      hence on $\kappa(v_E) \subseteq \kappa(v_\star)$, so $\sigma|_E \in I$
      and $\sigma \in I \cap N$.
\item Conversely, if $\sigma \in I \cap N$ then $\sigma$ fixes $K_{C'}$
      (whence acts trivially on $\kappa(V_{C'})$) and $\sigma|_E$ acts
      trivially on $\kappa(v_E)$.  Under separability of
      $\kappa(v_E)/\kappa(V_F)$ one has
      $\kappa(v_\star) = \kappa(v_E) \cdot \kappa(V_{C'})$, so $\sigma$
      acts trivially on $\kappa(v_\star)$, i.e.\ $\sigma \in I'$.
\end{itemize}
Hence $I \cap N = 1$ inside $H$.

\paragraph{Step 4: Cyclic prime-power structure forces $I = 1$.}
By (H3) the group $H \leq G = \Z/p^r\Z$ is cyclic of prime-power order, so
its subgroup lattice is a chain.  Consequently, two subgroups
$A, B \leq H$ satisfy $A \cap B = 1$ iff $A = 1$ or $B = 1$.  Applied to
$I \cap N = 1$ with $N \neq 1$, this gives $I = 1$.

Thus $E/F$ is unramified at $V_F$.  Combined with the unramified extension
$V_F/V_C$ of Step~1, $E/K_C$ is unramified at $V_C$, i.e.\ $\cP^{[d]_G}\to
U^{(d)}$ is unramified at $\eta_C$.
\end{proof}

\begin{rem*}\label{rem:where-each-hypothesis-bites}
The four hypotheses now play disjoint roles:
\begin{itemize}
  \item (H1) is used in Step~1 to produce the unramified extension
        $V_F/V_C$ and to close the tower at the end.
  \item (H2) is used in Step~3 to give $I' = 1$.
  \item (H3) is used \emph{only} in Step~4, to convert
        ``$I \cap N = 1$ with $N \neq 1$'' into $I = 1$.  For
        $G = (\Z/2)^2$ two nontrivial subgroups can intersect trivially,
        and the conclusion can fail.
  \item (H4) is used \emph{only} in Step~3, to validate the
        inertia-in-compositum identity \eqref{inertia-intersection}.
        It is automatic in characteristic $0$ and in the tame case.
\end{itemize}
The bounded-ramification assumption on $\cP$ at $p_\infty$ is upstream,
ensuring that $f_X \colon X_{C'} \dashrightarrow X_C$ sends
$\eta_{C'} \mapsto \eta_C$ and thus produces the finite extension of DVRs
$V_C \hookrightarrow V_{C'}$ on which Steps~1 and~3 rely.
\end{rem*}

% #############################################################################
\section{A second revised proof: dropping \texorpdfstring{(H4)}{(H4)}}\label{sec:main-revised-2}
% #############################################################################

The separability hypothesis (H4) introduced in \Cref{thm:main-revised} is
\emph{not} needed: separability of the residue extensions in question is
automatic for any finite Galois extension of henselian DVRs, regardless of
the residue characteristic or perfectness of the base.  The argument is a
clean application of the henselian fundamental equality $|G| = e \cdot f$
together with the bound $|\mathrm{Aut}_F(E)| \leq [E:F]_s$.  We restate the
theorem without (H4) and replace the relevant paragraph of Step~3
accordingly; everything else carries over verbatim from
\Cref{sec:main-revised}.

\begin{theorem}\label{thm:main-revised-2}
Assume
\begin{enumerate}
  \item[\textnormal{(H1)}] $(\cP_{G/H})^{[d]} \to U^{(d)}$ is unramified at
        $\eta_C$;
  \item[\textnormal{(H2)}] $\cP^{[d]_H} \to (\cP_{G/H})^{(d)}$ is unramified at
        $\eta_{C'}$;
  \item[\textnormal{(H3)}] $G = \Z/p^r\Z$ is cyclic of prime-power order.
\end{enumerate}
Then $\cP^{[d]_G} \to U^{(d)}$ is unramified at $\eta_C$.
\end{theorem}

The proof is identical to that of \Cref{thm:main-revised} except for the
verification of the inertia-in-compositum identity
\eqref{eq:inertia-intersection} in Step~3.  We record only this revised
fragment.

\begin{proof}[Revised Step~3 (replacement)]
With notation as in the proof of \Cref{thm:main-revised}, set
$M := E \cap K_{C'}$ and
$N := \Gal(E\cdot K_{C'}/K_{C'}) \cong \Gal(E/M) \leq H$, with $N \neq 1$
by Step~2.  Hypothesis (H2) gives, as before,
\[
I' \;:=\; I_{v_\star}\!\bigl(E\cdot K_{C'}/K_{C'}\bigr) \;=\; 1 ,
\]
and we claim
\begin{equation}\label{eq:inertia-intersection-2}
I' \;=\; I \cap N \qquad \text{inside } H ,
\end{equation}
where $I = I_{v_E}(E/F)$.

\paragraph{Separability of residue extensions is automatic.}
We first show, without any hypothesis on $k$ or on the characteristic, that
both residue extensions $\kappa(v_E)/\kappa(V_F)$ and
$\kappa(V_{C'})/\kappa(V_F)$ are separable.

Let $L/K$ be any finite Galois extension of henselian DVRs with Galois group
$\Gamma$, and write $e := |I_\Gamma|$ for the ramification index,
$f := [\kappa_L : \kappa_K]$ for the (total) residue degree.  Henselicity
gives a unique extension of the place of $K$ to $L$, and the fundamental
equality reads
\[
|\Gamma| \;=\; e \cdot f
\qquad\text{(Serre, \emph{Local Fields}, Ch.~I, \S7, Prop.~22).}
\]
Since $I_\Gamma$ is the kernel of the natural action
$\Gamma \to \mathrm{Aut}_{\kappa_K}(\kappa_L)$, the quotient $\Gamma/I_\Gamma$
embeds in $\mathrm{Aut}_{\kappa_K}(\kappa_L)$.  Combining with the standard
bound $|\mathrm{Aut}_{\kappa_K}(\kappa_L)| \leq [\kappa_L : \kappa_K]_s$
(where $[\,\cdot\,]_s$ denotes the separable degree) yields
\[
f \;=\; |\Gamma/I_\Gamma|
   \;\leq\; \bigl|\mathrm{Aut}_{\kappa_K}(\kappa_L)\bigr|
   \;\leq\; [\kappa_L : \kappa_K]_s
   \;\leq\; f .
\]
All inequalities collapse, forcing $f = [\kappa_L : \kappa_K]_s$, i.e.\ the
residue extension $\kappa_L/\kappa_K$ is separable.  No assumption on
$\mathrm{char}\,k$ or perfectness of $\kappa_K$ was used.

Applying this to $E/F$ at the place $v_E$ (after passing to the
henselisation $V_F^h$, where the decomposition group becomes the full local
Galois group) shows $\kappa(v_E)/\kappa(V_F)$ is separable.  Applying it
analogously to $K_{C'}/F$ at the place $V_{C'}$ shows
$\kappa(V_{C'})/\kappa(V_F)$ is separable.

\paragraph{Inertia identity.}
With both residue extensions separable, the standard compositum identity
\[
\kappa(v_\star) \;=\; \kappa(v_E) \cdot \kappa(V_{C'})
\]
holds (Bourbaki, \emph{Algèbre Commutative}, Ch.~VI, \S8, n.\textsuperscript{o}~3).
Both inclusions of \eqref{eq:inertia-intersection-2} then follow:
\begin{itemize}
\item If $\sigma' \in I'$ then $\sigma'$ acts trivially on $\kappa(v_\star)$,
      hence on $\kappa(v_E) \subseteq \kappa(v_\star)$, so $\sigma'|_E \in I$
      and $\sigma' \in I \cap N$.
\item If $\sigma' \in I \cap N$ then $\sigma'$ fixes $K_{C'}$ (whence acts
      trivially on $\kappa(V_{C'})$) and $\sigma'|_E$ acts trivially on
      $\kappa(v_E)$.  Since $\kappa(v_\star) = \kappa(v_E)\cdot\kappa(V_{C'})$,
      $\sigma'$ acts trivially on $\kappa(v_\star)$, i.e.\ $\sigma' \in I'$.
\end{itemize}
Hence $I \cap N = 1$ inside $H$.  Steps~1, 2, 4 of
\Cref{thm:main-revised} now conclude the proof of
\Cref{thm:main-revised-2}.
\end{proof}

\begin{rem*}\label{rem:why-h4-was-unnecessary}
The henselian fundamental equality $|G| = e \cdot f$ is doing the work that
(H4) was supposed to do.  The faithfulness inclusion
$\Gamma/I_\Gamma \hookrightarrow \mathrm{Aut}_{\kappa_K}(\kappa_L)$ on its
own would only show $|\Gamma/I_\Gamma| \leq [\kappa_L:\kappa_K]_s$, which is
a separability statement \emph{only when} we also know
$|\Gamma/I_\Gamma| = f$.  The latter is exactly the henselian equality, and
it is unconditional --- which is why no assumption on the base field or its
characteristic appears in \Cref{thm:main-revised-2}.

In particular, the residue fields $\kappa(\eta_C) = \kappa(p_\infty)(t_1,
\dots, t_{d-1})$ etc.\ that arise in our blowup setup are typically
imperfect in positive characteristic, yet the argument still works: the
relevant separability is of the \emph{Galois} residue extensions $\kappa(v_E)
/\kappa(V_F)$ and $\kappa(V_{C'})/\kappa(V_F)$, which are forced to be
separable irrespective of how imperfect the base $\kappa(V_F)$ is.

The earlier objection (that a direct symmetric-power-plus-blowup computation
of $\kappa(\eta_{C'})/\kappa(\eta_C)$ in the wild Artin--Schreier case seems
to produce a purely inseparable Frobenius-like map) misread the leading
term: in the relevant blowup grading, the correction terms in the local
expansion $\pi = (\pi')^p + (\pi')^{2p-1} + \cdots$ are of \emph{lower}
blowup-valuation than the naive $p$-th-power term, and they dominate.  The
extension $V_{C'}/V_C$ is in fact unramified in that example, consistent
with the Galois argument above.
\end{rem*}


\printbibliography

\end{document}